%============================================================
% Rapport Technique – Projet de Compilation TPC → NASM
% Vincent PLESSY – Licence Informatique 2025-2026
%============================================================
\documentclass[12pt,a4paper]{article}

%--- Encodage & langue ---
\usepackage[utf8]{inputenc}
\usepackage[T1]{fontenc}
\usepackage[french]{babel}

%--- Mise en page ---
\usepackage[top=2.5cm, bottom=2.5cm, left=2.5cm, right=2.5cm]{geometry}
\usepackage{fancyhdr}
\usepackage{titlesec}
\usepackage{parskip}

%--- Couleurs ---
\usepackage{xcolor}
\definecolor{bleutpc}{RGB}{0,70,170}
\definecolor{gristpc}{RGB}{100,100,100}
\definecolor{vertcomment}{RGB}{0,128,0}

%--- Listings (code) ---
\usepackage{listings}
\lstset{
  basicstyle=\ttfamily\small,
  keywordstyle=\color{bleutpc}\bfseries,
  commentstyle=\color{vertcomment}\itshape,
  stringstyle=\color{orange},
  numberstyle=\tiny\color{gristpc},
  numbers=left,
  numbersep=8pt,
  frame=single,
  rulecolor=\color{gristpc!50},
  backgroundcolor=\color{gray!5},
  breaklines=true,
  breakatwhitespace=true,
  tabsize=4,
  showstringspaces=false,
  captionpos=b,
}
\lstdefinelanguage{TPC}{
  morekeywords={int,char,void,struct,if,else,while,return},
  morecomment=[l]{//},
  morecomment=[s]{/*}{*/},
  morestring=[b]",
  morestring=[b]',
  sensitive=true,
}

%--- Tables ---
\usepackage{booktabs}
\usepackage{array}
\usepackage{tabularx}
\usepackage{longtable}

%--- Figures & flottants ---
\usepackage{float}
\usepackage{graphicx}

%--- Liens hypertexte ---
\usepackage[colorlinks=true, linkcolor=bleutpc, urlcolor=bleutpc, citecolor=bleutpc]{hyperref}

%--- En-têtes / pieds de page ---
\pagestyle{fancy}
\fancyhf{}
\fancyhead[L]{\nouppercase{\leftmark}}
\fancyhead[R]{\textcolor{gristpc}{Vincent PLESSY}}
\fancyfoot[C]{\thepage}
\renewcommand{\headrulewidth}{0.4pt}
\renewcommand{\footrulewidth}{0pt}

%--- Titres colorés ---
\titleformat{\section}
  {\color{bleutpc}\large\bfseries}{\thesection}{1em}{}[\color{bleutpc}\hrule]
\titleformat{\subsection}
  {\color{bleutpc}\normalsize\bfseries}{\thesubsection}{1em}{}
\titleformat{\subsubsection}
  {\color{bleutpc}\normalsize\itshape}{\thesubsubsection}{1em}{}

%--- Commandes utiles ---
\newcommand{\tpc}{\texttt{TPC}}
\newcommand{\tpcc}{\texttt{tpcc}}
\newcommand{\nasm}{\textsc{NASM}}

%============================================================
\begin{document}
%============================================================

%--- Page de titre ---
\begin{titlepage}
\centering
\vspace*{2cm}

{\color{bleutpc}\rule{\linewidth}{1.5pt}}\\[0.4cm]
{\Huge\bfseries\color{bleutpc} Rapport Technique}\\[0.3cm]
{\large\color{gristpc} Projet de Compilation}\\[0.1cm]
{\color{bleutpc}\rule{\linewidth}{1.5pt}}

\vspace{1.5cm}

{\LARGE\bfseries Compilateur pour le Langage TPC}\\[0.5cm]
{\large\textit{De l'analyse syntaxique à la génération de code NASM}}

\vspace{2cm}

{\large\bfseries Vincent PLESSY}\\[0.3cm]
{\normalsize Licence d'Informatique – Année 2025-2026}\\[0.2cm]
{\normalsize Mars 2026}

\vfill

{\small\color{gristpc}
  Université – Licence Informatique\\
  Projet de Compilation – Semestre 6\\
  Langage : TPC (sous-ensemble du C)\\
  Cible : NASM ELF64 (Linux AMD64)
}

\vspace{1cm}
{\color{bleutpc}\rule{\linewidth}{0.4pt}}
\end{titlepage}

%--- Table des matières ---
\tableofcontents
\newpage

%============================================================
\section{Introduction}
%============================================================

Ce rapport documente la réalisation du compilateur \tpcc{} pour le langage \tpc{}
(\textit{Tiny Procedural C}), dans le cadre du projet de Compilation de Licence~3.

Ce projet est la continuation directe du projet d'Analyse Syntaxique (Semestre~5),
qui avait produit l'analyseur \texttt{tpcas} capable de détecter les erreurs
lexicales et syntaxiques et de construire un Arbre de Syntaxe Abstraite (ASA).
Le présent compilateur étend ce travail avec :

\begin{itemize}
  \item une \textbf{analyse sémantique} complète (tables des symboles, vérification
        des types, contrôle des déclarations) ;
  \item une \textbf{génération de code} vers l'assembleur \nasm{} (format ELF64,
        architecture AMD64 Linux) ;
  \item une suite de tests automatisée couvrant les quatre catégories attendues.
\end{itemize}

L'ensemble du code source peut être compilé avec un simple \texttt{make}.

%============================================================
\section{Spécifications Techniques}
%============================================================

\subsection{Environnement de développement}

\begin{table}[H]
\centering
\begin{tabularx}{\linewidth}{lX}
\toprule
\textbf{Composant} & \textbf{Version / Description} \\
\midrule
Système d'exploitation & Linux AMD64 (Debian/Ubuntu) \\
Compilateur C          & GCC 11+ \\
Analyseur lexical      & Flex 2.6+ \\
Générateur de parseur  & Bison 3.8+ \\
Assembleur cible       & NASM 2.15+ (format ELF64) \\
Éditeur de liens       & ld (GNU Binutils) \\
Outil de construction  & GNU Make \\
\bottomrule
\end{tabularx}
\caption{Outils utilisés}
\end{table}

\subsection{Structure du projet}

\begin{verbatim}
project/
├── makefile
├── run_tests.sh         # Script de tests automatiques
├── bin/
│   └── tpcc             # Binaire compilé
├── src/
│   ├── tpcc.lex         # Analyseur lexical (Flex)
│   ├── tpcc.y           # Parseur + point d'entrée (Bison)
│   ├── tree.h / tree.c  # Module ASA
│   ├── symtable.h / symtable.c   # Tables des symboles
│   ├── semantic.h / semantic.c   # Analyse sémantique
│   └── codegen.h / codegen.c     # Génération de code NASM
├── test/
│   ├── good/            # 12 programmes corrects
│   ├── syn-err/         # 5 erreurs lexicales/syntaxiques
│   ├── sem-err/         # 10 erreurs sémantiques
│   └── warn/            # 3 programmes avec avertissements
└── rep/
    └── RAPPORT.tex      # Ce document
\end{verbatim}

\subsection{Métriques du code source}

\begin{table}[H]
\centering
\begin{tabular}{lrrr}
\toprule
\textbf{Fichier} & \textbf{Lignes} & \textbf{Rôle} \\
\midrule
\texttt{tpcc.lex}      & $\sim$110 & Lexer Flex \\
\texttt{tpcc.y}        & $\sim$250 & Parser Bison + \texttt{main()} \\
\texttt{tree.c}        & $\sim$120 & Module ASA \\
\texttt{symtable.c}    & $\sim$220 & Tables des symboles \\
\texttt{semantic.c}    & $\sim$380 & Analyse sémantique \\
\texttt{codegen.c}     & $\sim$580 & Génération NASM \\
\midrule
\textbf{Total}         & $\sim$1660 & \\
\bottomrule
\end{tabular}
\caption{Taille des fichiers sources}
\end{table}

\subsection{Codes de retour}

\begin{table}[H]
\centering
\begin{tabular}{cl}
\toprule
\textbf{Code} & \textbf{Signification} \\
\midrule
0 & Compilation réussie (fichier \texttt{.asm} produit) \\
1 & Erreur lexicale ou syntaxique \\
2 & Erreur sémantique \\
3 & Erreur interne / fichier introuvable \\
\bottomrule
\end{tabular}
\caption{Codes de retour de \tpcc{}}
\end{table}

%============================================================
\section{Analyse Lexicale}
%============================================================

\subsection{Tokens reconnus}

L'analyseur lexical \texttt{tpcc.lex} reconnaît les catégories suivantes :

\begin{table}[H]
\centering
\begin{tabularx}{\linewidth}{lX}
\toprule
\textbf{Catégorie} & \textbf{Description} \\
\midrule
\texttt{TYPE}       & \texttt{int}, \texttt{char}, \texttt{void} \\
\texttt{STRUCT}     & Mot-clé \texttt{struct} \\
\texttt{IF / ELSE}  & Mots-clés de condition \\
\texttt{WHILE}      & Boucle \\
\texttt{RETURN}     & Instruction de retour \\
\texttt{IDENT}      & Identificateurs (\texttt{[a-zA-Z\_][a-zA-Z0-9\_]*}) \\
\texttt{NUM}        & Entiers (\texttt{[0-9]+}) \\
\texttt{CHARACTER}  & Littéral caractère (\texttt{'x'}, \texttt{'\textbackslash n'}, etc.) \\
\texttt{ADDSUB}     & \texttt{+}, \texttt{-} \\
\texttt{DIVSTAR}    & \texttt{*}, \texttt{/}, \texttt{\%} \\
\texttt{EQ}         & \texttt{==}, \texttt{!=} \\
\texttt{ORDER}      & \texttt{<}, \texttt{<=}, \texttt{>}, \texttt{>=} \\
\texttt{AND / OR}   & \texttt{\&\&}, \texttt{||} \\
\texttt{NOT}        & \texttt{!} \\
Ponctuations        & \texttt{( ) \{ \} [ ] ; , .} \\
\bottomrule
\end{tabularx}
\caption{Tokens du lexer}
\end{table}

\subsection{Gestion des erreurs lexicales}

Tout caractère non reconnu provoque l'émission d'un message d'erreur sur
\texttt{stderr} et la valeur \texttt{lexical\_error} est positionnée afin de
retourner le code~1 en fin de parsing.

%============================================================
\section{Analyse Syntaxique}
%============================================================

\subsection{Grammaire TPC}

Le parseur \texttt{tpcc.y} implémente la grammaire complète du langage TPC.
Les productions principales sont :

\begin{lstlisting}[language=C, caption={Extrait de la grammaire Bison}]
Prog      : DeclList FunctList
          | FunctList
          ;

DeclList  : DeclList Decl
          | Decl
          ;

Decl      : TYPE DeclVars ';'
          | TYPE IDENT '[' NUM ']' ';'
          | StructDef ';'
          ;

FunctList : FunctList Funct
          | Funct
          ;

Funct     : TypeFunct IDENT '(' Params ')' Body
          ;

Instr     : Assign ';'
          | StructAff ';'
          | If
          | While
          | Return ';'
          | FunctCall ';'
          | '{' Body '}'
          | ';'
          ;
\end{lstlisting}

\subsection{Priorités des opérateurs}

Les priorités sont déclarées via les directives Bison \texttt{\%left} /
\texttt{\%right}, du plus faible au plus fort :

\begin{enumerate}
  \item \texttt{||} (ou logique)
  \item \texttt{\&\&} (et logique)
  \item \texttt{==}, \texttt{!=} (égalité)
  \item \texttt{<}, \texttt{<=}, \texttt{>}, \texttt{>=} (ordre)
  \item \texttt{+}, \texttt{-} (addition/soustraction)
  \item \texttt{*}, \texttt{/}, \texttt{\%} (multiplication/division)
  \item \texttt{!}, unaire \texttt{-} (opérateurs unaires)
\end{enumerate}

%============================================================
\section{Arbre de Syntaxe Abstraite (ASA)}
%============================================================

\subsection{Structure des nœuds}

Le module \texttt{tree.h} / \texttt{tree.c} est identique à celui du projet~S1.
Chaque nœud de l'ASA est représenté par la structure :

\begin{lstlisting}[language=C, caption={Structure d'un nœud ASA}]
typedef struct Node {
    label_t      label;       /* type de nœud */
    struct Node *firstChild;
    struct Node *nextSibling;
    int          lineno;
    char         ident[MAXNAME];  /* pour Ident, Character */
    int          num;             /* pour Num */
    int          type;            /* type TPC (TY_INT, etc.) */
} Node;
\end{lstlisting}

\subsection{Labels de l'ASA}

\begin{table}[H]
\centering
\begin{tabularx}{\linewidth}{lX}
\toprule
\textbf{Label} & \textbf{Description} \\
\midrule
\texttt{Prog}          & Racine du programme \\
\texttt{DeclVars}      & Liste de déclarations de variables \\
\texttt{Declarator}    & Un déclarant (nom + éventuel tableau) \\
\texttt{Funct}         & Définition de fonction \\
\texttt{Params}        & Liste de paramètres formels \\
\texttt{Param}         & Un paramètre \\
\texttt{Body}          & Corps de fonction (liste d'instructions) \\
\texttt{InstrAssign}   & Affectation \\
\texttt{InstrIf}       & Condition simple \\
\texttt{InstrIfElse}   & Condition avec alternative \\
\texttt{InstrWhile}    & Boucle \texttt{while} \\
\texttt{InstrCall}     & Appel de procédure \\
\texttt{InstrReturn}   & \texttt{return expr} \\
\texttt{InstrReturnVoid} & \texttt{return} sans valeur \\
\texttt{InstrBlock}    & Bloc \texttt{\{...\}} \\
\texttt{F}             & Appel de fonction (dans expression) \\
\texttt{Ident}         & Identificateur \\
\texttt{Num}           & Entier littéral \\
\texttt{Character}     & Caractère littéral \\
\texttt{OpAddsub}      & Opérateur \texttt{+} / \texttt{-} \\
\texttt{OpDivstar}     & Opérateur \texttt{*} / \texttt{/} / \texttt{\%} \\
\texttt{OpEq}          & Opérateur \texttt{==} / \texttt{!=} \\
\texttt{OpOrder}       & Opérateurs de comparaison \\
\texttt{OpAnd / OpOr}  & Opérateurs logiques \\
\texttt{OpNot}         & Négation logique \\
\texttt{OpUnaryMinus}  & Moins unaire \\
\texttt{StructDef / StructVar / StructAff} & Structures \\
\bottomrule
\end{tabularx}
\caption{Principaux labels de l'ASA}
\end{table}

%============================================================
\section{Tables des Symboles}
%============================================================

\subsection{Architecture}

Le module \texttt{symtable.h} / \texttt{symtable.c} gère quatre tables globales :

\begin{description}
  \item[\texttt{g\_structs}] Définitions de structures (\texttt{struct}).
  \item[\texttt{g\_vars}]    Variables globales.
  \item[\texttt{g\_funcs}]   Fonctions déclarées (incluant les fonctions prédéfinies).
  \item[\texttt{l\_vars}]    Variables locales et paramètres de la fonction courante.
\end{description}

\subsection{Structure d'un symbole}

\begin{lstlisting}[language=C, caption={Structure Symbol}]
typedef struct Symbol {
    char     name[MAX_NAME];
    SymKind  kind;        /* SYM_GVAR, SYM_LVAR, SYM_PARAM, SYM_FUNC, SYM_SDEF */
    TpcType  type;        /* {kind: TY_INT/TY_CHAR/TY_VOID/TY_STRUCT, struct_name} */
    int      offset;      /* offset rbp pour locaux/params (négatif) */
    int      size;        /* taille en octets */
    int      nparams;     /* nb de paramètres (pour les fonctions) */
    Symbol  *params;      /* liste chaînée des paramètres */
    Field    fields[MAX_FIELDS];
    int      nfields;     /* nb de champs (pour les structs) */
    int      struct_size; /* taille totale de la struct */
    struct Symbol *next;
} Symbol;
\end{lstlisting}

\subsection{Allocation des variables locales}

Chaque variable locale ou paramètre occupe \textbf{8~octets} sur la pile
(indépendamment de son type effectif), ce qui simplifie le calcul des offsets.
L'offset de la première variable est $-8$ par rapport à \texttt{rbp}, puis
$-16$, $-24$, etc.

\begin{lstlisting}[language=C, caption={Allocation d'une variable locale}]
void sym_add_lvar(const char *name, TpcType type) {
    Symbol *s = calloc(1, sizeof(*s));
    strncpy(s->name, name, MAX_NAME-1);
    s->kind   = SYM_LVAR;
    s->type   = type;
    s->offset = next_offset;   /* négatif, commence à -8 */
    s->size   = 8;
    next_offset -= 8;
    /* insertion en tête de l_vars */
    s->next = l_vars;
    l_vars  = s;
}
\end{lstlisting}

\subsection{Fonctions prédéfinies}

Les quatre fonctions d'E/S sont pré-enregistrées dans \texttt{g\_funcs}
lors de l'initialisation :

\begin{table}[H]
\centering
\begin{tabular}{lll}
\toprule
\textbf{Nom} & \textbf{Paramètre} & \textbf{Retour} \\
\midrule
\texttt{putchar} & \texttt{int c}  & \texttt{void} \\
\texttt{putint}  & \texttt{int n}  & \texttt{void} \\
\texttt{getchar} & \textit{aucun}  & \texttt{int} \\
\texttt{getint}  & \textit{aucun}  & \texttt{int} \\
\bottomrule
\end{tabular}
\caption{Fonctions prédéfinies TPC}
\end{table}

%============================================================
\section{Analyse Sémantique}
%============================================================

\subsection{Vérifications effectuées}

L'analyse sémantique (\texttt{semantic.c}) parcourt l'ASA après le parsing
et effectue les vérifications suivantes :

\begin{table}[H]
\centering
\begin{tabularx}{\linewidth}{lX}
\toprule
\textbf{Vérification} & \textbf{Code retour} \\
\midrule
Variable non déclarée                     & erreur (2) \\
Variable déclarée deux fois               & erreur (2) \\
Fonction non déclarée                     & erreur (2) \\
Fonction déclarée deux fois               & erreur (2) \\
Nombre d'arguments incorrect              & erreur (2) \\
Utilisation d'une fonction \texttt{void} comme expression & erreur (2) \\
Absence de fonction \texttt{main}         & erreur (2) \\
\texttt{main} avec mauvaise signature     & erreur (2) \\
Redéclaration locale d'un paramètre       & erreur (2) \\
Affectation \texttt{int} → \texttt{char}  & avertissement \\
\texttt{return int} dans fonction \texttt{char} & avertissement \\
\bottomrule
\end{tabularx}
\caption{Vérifications sémantiques}
\end{table}

\subsection{Algorithme général}

\begin{enumerate}
  \item Pré-enregistrement des fonctions prédéfinies.
  \item Parcours des définitions de structures (\texttt{StructDef}).
  \item Enregistrement des variables globales (\texttt{DeclVars}).
  \item Pour chaque fonction :
    \begin{enumerate}
      \item Enregistrement dans \texttt{g\_funcs}.
      \item Réinitialisation de \texttt{l\_vars}.
      \item Enregistrement des paramètres formels.
      \item Parcours récursif du corps (\texttt{Body}).
    \end{enumerate}
  \item Vérification de la présence de \texttt{main()}.
\end{enumerate}

\subsection{Système de types}

La compatibilité de types suit les règles :

\begin{itemize}
  \item \texttt{int} → \texttt{int} : compatible (OK)
  \item \texttt{char} → \texttt{char} : compatible (OK)
  \item \texttt{char} → \texttt{int} : compatible implicite (OK, pas d'avertissement)
  \item \texttt{int} → \texttt{char} : compatible avec avertissement (code~0, message sur \texttt{stderr})
  \item Toute autre combinaison : incompatible (erreur, code~2)
\end{itemize}

%============================================================
\section{Génération de Code NASM}
%============================================================

\subsection{Architecture générale}

Le module \texttt{codegen.c} traduit l'ASA en assembleur \nasm{} ELF64.
Le fichier généré est structuré en trois sections :

\begin{description}
  \item[\texttt{section .bss}] Variables globales (réservation de 8~octets par variable).
  \item[\texttt{section .text}] Code : \texttt{\_start}, toutes les fonctions, puis les fonctions prédéfinies.
\end{description}

\begin{lstlisting}[language={[x86masm]Assembler}, caption={Structure d'un fichier .asm généré}]
section .bss
    _x:  resb 8          ; variable globale x

section .text
    global _start
    extern putchar, putint, getchar, getint  ; si non générés

_start:
    call main
    mov  rdi, rax
    mov  rax, 60
    syscall

main:
    push rbp
    mov  rbp, rsp
    sub  rsp, 16
    ; corps de la fonction
    xor  eax, eax
    leave
    ret
\end{lstlisting}

\subsection{Convention d'appel AMD64}

Le compilateur suit la convention \textit{System V AMD64 ABI} :

\begin{table}[H]
\centering
\begin{tabular}{cl}
\toprule
\textbf{Registre} & \textbf{Rôle} \\
\midrule
\texttt{rdi} & 1\textsuperscript{er} argument entier \\
\texttt{rsi} & 2\textsuperscript{e} argument entier \\
\texttt{rdx} & 3\textsuperscript{e} argument entier \\
\texttt{rcx} & 4\textsuperscript{e} argument entier \\
\texttt{r8}  & 5\textsuperscript{e} argument entier \\
\texttt{r9}  & 6\textsuperscript{e} argument entier \\
\texttt{rax} & Valeur de retour \\
\bottomrule
\end{tabular}
\caption{Registres AMD64 utilisés}
\end{table}

\subsection{Prologue et épilogue de fonction}

Chaque fonction suit le schéma :

\begin{lstlisting}[language={[x86masm]Assembler}, caption={Prologue/épilogue standard}]
nom_fonction:
    push rbp
    mov  rbp, rsp
    sub  rsp, N        ; N = taille pile locale (alignée 16 octets)
    ; sauvegarde des paramètres :
    mov  [rbp-8],  rdi
    mov  [rbp-16], rsi
    ; ...
    ; corps
    xor  eax, eax      ; valeur de retour par défaut (0)
    leave
    ret
\end{lstlisting}

L'alignement de la pile respecte la contrainte AMD64 : après \texttt{push rbp},
\texttt{rsp} est décalé de~8. On alloue donc $N$ tel que $N \equiv 8 \pmod{16}$
(garantit que \texttt{rsp} est multiple de~16 avant tout \texttt{call}).

\subsection{Génération des expressions}

La fonction \texttt{cg\_expr()} place le résultat dans \texttt{eax} :

\begin{table}[H]
\centering
\begin{tabularx}{\linewidth}{lX}
\toprule
\textbf{Nœud ASA} & \textbf{Code généré} \\
\midrule
\texttt{Num n}             & \texttt{mov eax, n} \\
\texttt{Character 'c'}     & \texttt{mov eax, ord(c)} \\
\texttt{Ident x} (local)   & \texttt{mov eax, [rbp+offset]} \\
\texttt{Ident x} (global)  & \texttt{mov eax, [\_x]} \\
\texttt{OpAddsub +}        & éval gauche, push, éval droite, pop rcx, add \\
\texttt{OpDivstar *}       & éval gauche, push, éval droite, pop rcx, imul \\
\texttt{OpDivstar /}       & idiv (quotient dans eax) \\
\texttt{OpEq ==}           & cmp + sete al + movzx \\
\texttt{OpAnd}             & court-circuit : saute si gauche = 0 \\
\texttt{OpOr}              & court-circuit : saute si gauche $\neq$ 0 \\
\texttt{OpNot}             & test eax,eax + sete al \\
\texttt{OpUnaryMinus}      & neg eax \\
\texttt{F (appel)}         & push args, call, (pop args si besoin) \\
\bottomrule
\end{tabularx}
\caption{Génération des expressions}
\end{table}

\subsection{Génération des instructions}

\begin{description}
  \item[\texttt{InstrAssign}] Évalue l'expression droite dans \texttt{eax}, stocke
        dans la variable gauche (\texttt{[rbp+offset]} ou \texttt{[\_x]}).
  \item[\texttt{InstrIf}] Évalue la condition, saute à \texttt{\_L<n>\_end} si faux.
  \item[\texttt{InstrIfElse}] Saute à \texttt{\_L<n>\_else} si faux, puis à
        \texttt{\_L<n>\_end} en fin de branche vraie.
  \item[\texttt{InstrWhile}] Étiquette \texttt{\_L<n>\_start}, teste, corps,
        saute à \texttt{start}, étiquette \texttt{\_L<n>\_end}.
  \item[\texttt{InstrReturn}] Évalue l'expression dans \texttt{eax}, puis
        \texttt{leave} / \texttt{ret}.
  \item[\texttt{InstrBlock}] Parcours récursif des instructions enfants.
\end{description}

\subsection{Fonctions d'E/S prédéfinies}

Les quatre fonctions d'E/S sont implémentées directement en assembleur via les
appels système Linux (sans bibliothèque C) :

\begin{lstlisting}[language={[x86masm]Assembler}, caption={Implémentation de putchar en NASM}]
putchar:
    push rbp
    mov  rbp, rsp
    sub  rsp, 16
    mov  [rbp-8], rdi       ; sauvegarde argument
    mov  byte [rbp-9], dil  ; octet à afficher
    lea  rsi, [rbp-9]       ; pointeur vers le caractère
    mov  rax, 1             ; syscall write
    mov  rdi, 1             ; fd = stdout
    mov  rdx, 1             ; longueur = 1
    syscall
    leave
    ret
\end{lstlisting}

\begin{lstlisting}[language={[x86masm]Assembler}, caption={Implémentation de putint en NASM (extrait)}]
putint:
    ; convertit l'entier en chaîne ASCII, puis syscall write
    ; gère le signe négatif, les chiffres 0-9
    ; utilise une pile de chiffres sur la pile locale
\end{lstlisting}

%============================================================
\section{Extension : Structures}
%============================================================

\subsection{Déclaration et utilisation}

Le langage TPC supporte les structures :

\begin{lstlisting}[language=TPC, caption={Exemple d'utilisation d'une structure}]
struct Point {
    int x;
    int y;
};

int main(void) {
    struct Point p;
    p.x = 10;
    p.y = 20;
    putint(p.x + p.y);
    return 0;
}
\end{lstlisting}

\subsection{Représentation en mémoire}

Les structures sont stockées dans la table \texttt{g\_structs}. Chaque champ
est représenté par un \texttt{Field} (nom + type + offset interne). Le
compilateur supporte l'accès aux champs via le nœud \texttt{StructAff} de l'ASA.

%============================================================
\section{Tests et Validation}
%============================================================

\subsection{Organisation des tests}

La suite de tests comprend \textbf{30~fichiers} répartis en quatre catégories :

\begin{table}[H]
\centering
\begin{tabular}{llr}
\toprule
\textbf{Répertoire} & \textbf{Catégorie} & \textbf{Nb fichiers} \\
\midrule
\texttt{test/good/}    & Programmes corrects (code 0)      & 12 \\
\texttt{test/syn-err/} & Erreurs lexicales/syntaxiques (code 1) & 5 \\
\texttt{test/sem-err/} & Erreurs sémantiques (code 2)      & 10 \\
\texttt{test/warn/}    & Programmes avec avertissements (code 0) & 3 \\
\midrule
\textbf{Total}         &                                   & \textbf{30} \\
\bottomrule
\end{tabular}
\caption{Répartition des tests}
\end{table}

\subsection{Détail des tests}

\begin{longtable}{lp{8cm}}
\toprule
\textbf{Fichier} & \textbf{Description} \\
\midrule
\endhead
\multicolumn{2}{l}{\textit{good/ – programmes corrects}} \\
\midrule
\texttt{01\_main\_simple.tpc}     & Programme minimal (\texttt{main} retourne 0) \\
\texttt{02\_variables\_int.tpc}   & Déclarations et affectations d'entiers \\
\texttt{03\_arithmetic.tpc}       & Opérations arithmétiques (+, -, *, /, \%) \\
\texttt{04\_if\_else.tpc}         & Conditionnelle avec branche \texttt{else} \\
\texttt{05\_while.tpc}            & Boucle \texttt{while} \\
\texttt{06\_functions.tpc}        & Définition et appel de fonctions \\
\texttt{07\_global\_vars.tpc}     & Variables globales \\
\texttt{08\_putint.tpc}           & Utilisation de \texttt{putint} \\
\texttt{09\_char\_vars.tpc}       & Variables de type \texttt{char} \\
\texttt{10\_struct.tpc}           & Déclaration et accès à une structure \\
\texttt{11\_recursion.tpc}        & Fonction récursive (factorielle) \\
\texttt{12\_logical\_ops.tpc}     & Opérateurs logiques (\&\&, ||, !) \\
\midrule
\multicolumn{2}{l}{\textit{syn-err/ – erreurs syntaxiques}} \\
\midrule
\texttt{e01\_missing\_semicolon.tpc} & Point-virgule manquant \\
\texttt{e02\_missing\_brace.tpc}     & Accolade fermante absente \\
\texttt{e03\_invalid\_expr.tpc}      & Expression invalide \\
\texttt{e04\_bad\_struct.tpc}        & Structure mal formée \\
\texttt{e05\_bad\_return.tpc}        & Opérateur sans opérande droite \\
\midrule
\multicolumn{2}{l}{\textit{sem-err/ – erreurs sémantiques}} \\
\midrule
\texttt{s01\_undeclared\_var.tpc}      & Variable non déclarée \\
\texttt{s02\_undeclared\_func.tpc}     & Fonction non déclarée \\
\texttt{s03\_duplicate\_global.tpc}    & Variable globale déclarée deux fois \\
\texttt{s04\_duplicate\_func.tpc}      & Fonction déclarée deux fois \\
\texttt{s05\_wrong\_args.tpc}          & Nombre d'arguments incorrect \\
\texttt{s06\_void\_as\_expr.tpc}       & Fonction \texttt{void} utilisée en expression \\
\texttt{s07\_no\_main.tpc}             & Absence de \texttt{main} \\
\texttt{s08\_bad\_main.tpc}            & \texttt{main} avec mauvaise signature \\
\texttt{s09\_type\_error.tpc}          & Incompatibilité de type \\
\texttt{s10\_local\_param\_same\_name.tpc} & Redéclaration locale d'un paramètre \\
\midrule
\multicolumn{2}{l}{\textit{warn/ – programmes avec avertissements}} \\
\midrule
\texttt{w01\_int\_to\_char.tpc}            & Affectation \texttt{int} → \texttt{char} \\
\texttt{w02\_int\_return\_in\_char\_func.tpc} & \texttt{return int} dans fonction \texttt{char} \\
\texttt{w03\_implicit\_char\_to\_int.tpc}  & Conversion implicite \texttt{char} → \texttt{int} (OK) \\
\bottomrule
\caption{Détail des fichiers de test}
\end{longtable}

\subsection{Résultats}

\begin{table}[H]
\centering
\begin{tabular}{lrrrr}
\toprule
\textbf{Catégorie} & \textbf{Total} & \textbf{PASS} & \textbf{FAIL} & \textbf{Taux} \\
\midrule
good    & 12 & 12 & 0 & 100\% \\
syn-err & 5  & 5  & 0 & 100\% \\
sem-err & 10 & 10 & 0 & 100\% \\
warn    & 3  & 3  & 0 & 100\% \\
\midrule
\textbf{Total} & \textbf{30} & \textbf{30} & \textbf{0} & \textbf{100\%} \\
\bottomrule
\end{tabular}
\caption{Résultats des tests automatiques}
\end{table}

\subsection{Script de tests}

Le script \texttt{run\_tests.sh} lance automatiquement les 30 tests et
affiche un rapport coloré. Il supporte l'option \texttt{--asm} pour
assembler et lier les programmes corrects avec \texttt{nasm} + \texttt{ld}.

\begin{verbatim}
$ ./run_tests.sh
==============================
 Tests du compilateur tpcc
==============================

--- Programmes corrects (code de retour attendu: 0) ---
  PASS [good] 01_main_simple.tpc (code=0)
  ...

--- Résultats finaux ---
 PASS  : 30 / 30
 FAIL  : 0 / 30
 Score : 100%
==============================
\end{verbatim}

%============================================================
\section{Difficultés Rencontrées}
%============================================================

\subsection{Alignement de la pile AMD64}

La convention AMD64 exige que \texttt{rsp} soit aligné sur 16~octets avant
toute instruction \texttt{call}. Après \texttt{push rbp}, \texttt{rsp} a déjà
perdu 8~octets d'alignement. La formule d'allocation utilisée est :

\[
N = \begin{cases}
  \text{total} & \text{si } \text{total} \equiv 8 \pmod{16} \\
  \left\lceil \frac{\text{total}+15}{16} \right\rceil \times 16 - 8 & \text{sinon}
\end{cases}
\]

où \texttt{total} est la somme des espaces nécessaires pour toutes les
variables locales et les paramètres.

\subsection{Gestion des codes de retour dans les tests}

Le script \texttt{run\_tests.sh} ne doit pas utiliser \texttt{set -e} (qui
provoquerait la sortie immédiate du script dès qu'une commande retourne un
code non nul). La capture du code de retour utilise la syntaxe :

\begin{lstlisting}[language=bash]
$COMPILER "$file" >/dev/null 2>&1; actual_ret=$?
\end{lstlisting}

\subsection{Tokens CHARACTER et séquences d'échappement}

Le token \texttt{CHARACTER} stocke le littéral sous forme de chaîne (type
\texttt{<str>} dans l'union Bison) pour pouvoir gérer les séquences
d'échappement (\texttt{'\textbackslash n'}, \texttt{'\textbackslash t'},
\texttt{'\textbackslash 0'}, etc.) lors de la génération de code.

\subsection{Variables locales vs paramètres}

Les paramètres formels sont insérés dans \texttt{l\_vars} \textit{en tête}
lors du remplissage de la table des symboles. À la génération, ils sont
parcourus dans l'ordre de la liste (tête en premier), qui correspond à
l'ordre inverse de déclaration. Le code de sauvegarde des registres AMD64
tient compte de cet ordre pour associer correctement chaque registre
(\texttt{rdi}, \texttt{rsi}, …) à son offset.

%============================================================
\section{Résultats et Performances}
%============================================================

\subsection{Bilan fonctionnel}

\begin{table}[H]
\centering
\begin{tabularx}{\linewidth}{Xl}
\toprule
\textbf{Fonctionnalité} & \textbf{Statut} \\
\midrule
Analyse lexicale complète          & \textcolor{green!60!black}{Implémenté} \\
Analyse syntaxique (grammaire TPC) & \textcolor{green!60!black}{Implémenté} \\
Construction de l'ASA              & \textcolor{green!60!black}{Implémenté} \\
Tables des symboles                & \textcolor{green!60!black}{Implémenté} \\
Analyse sémantique                 & \textcolor{green!60!black}{Implémenté} \\
Génération NASM (types int/char)   & \textcolor{green!60!black}{Implémenté} \\
Génération NASM (structures)       & \textcolor{green!60!black}{Implémenté} \\
Fonctions d'E/S (syscall)          & \textcolor{green!60!black}{Implémenté} \\
Gestion des avertissements         & \textcolor{green!60!black}{Implémenté} \\
Tests 30/30                        & \textcolor{green!60!black}{100\%} \\
\bottomrule
\end{tabularx}
\caption{État d'avancement du projet}
\end{table}

\subsection{Limitations connues}

\begin{itemize}
  \item Les tableaux sont reconnus syntaxiquement mais l'accès indicé
        (\texttt{a[i]}) n'est pas entièrement géré en génération de code.
  \item Les fonctions à plus de 6~paramètres ne sont pas supportées
        (limitation de la convention AMD64 sans passage par la pile).
  \item Pas de vérification des débordements arithmétiques.
\end{itemize}

%============================================================
\section{Conclusion}
%============================================================

Ce projet a permis de réaliser un compilateur complet pour le langage TPC,
couvrant toutes les étapes de la chaîne de compilation : analyse lexicale,
analyse syntaxique, construction de l'ASA, analyse sémantique et génération
de code assembleur NASM.

Les 30 tests automatiques passent tous avec succès (score 100\%), validant
le bon fonctionnement de l'ensemble. La génération de code respecte la
convention AMD64 et produit des exécutables ELF64 fonctionnels sur Linux.

Ce travail constitue une base solide pour des extensions futures : optimisation
du code généré, support complet des tableaux, ou encore gestion d'appels de
fonctions avec plus de 6~arguments.

%============================================================
\appendix
\section{Annexes}
%============================================================

\subsection{Exemple de programme TPC compilé}

\begin{lstlisting}[language=TPC, caption={Calcul de la factorielle (test/good/11\_recursion.tpc)}]
int fact(int n) {
    if (n <= 1) return 1;
    return n * fact(n - 1);
}

int main(void) {
    putint(fact(5));
    putchar('\n');
    return 0;
}
\end{lstlisting}

\subsection{Assembleur NASM généré (extrait)}

\begin{lstlisting}[language={[x86masm]Assembler}, caption={Code NASM pour fact(5)}]
fact:
    push rbp
    mov  rbp, rsp
    sub  rsp, 8
    mov  [rbp-8], rdi     ; paramètre n

    ; if (n <= 1)
    mov  eax, [rbp-8]
    push rax
    mov  eax, 1
    pop  rcx
    cmp  ecx, eax
    setle al
    movzx eax, al
    test eax, eax
    jz   _L1_end
    mov  eax, 1           ; return 1
    leave
    ret
_L1_end:
    ; return n * fact(n-1)
    mov  eax, [rbp-8]
    push rax
    mov  eax, [rbp-8]
    push rax
    mov  eax, 1
    pop  rcx
    sub  ecx, eax
    mov  eax, ecx
    mov  rdi, rax
    call fact
    pop  rcx
    imul ecx, eax
    mov  eax, ecx
    leave
    ret
\end{lstlisting}

\subsection{Compilation et utilisation}

\begin{verbatim}
# Compiler le projet
make

# Compiler un programme TPC
./bin/tpcc mon_prog.tpc       # produit mon_prog.asm

# Assembler et exécuter
nasm -f elf64 mon_prog.asm -o mon_prog.o
ld -o mon_prog mon_prog.o
./mon_prog

# Lancer tous les tests
./run_tests.sh

\end{verbatim}

\end{document}